\documentclass[letterpaper, 12pt,oneside]{article}
\usepackage{amsmath}
\usepackage{graphicx}
\usepackage{xcolor}
\graphicspath{{Imagenes/}}
\usepackage[utf8]{inputenc}
\usepackage{listings}
\usepackage[hidelinks]{hyperref}

\title{\Huge Taller de Herramientas Computacionales}
\author{Josué Artemio Hernández Rodríguez}
\date{24/Enero/2019}

\begin{document}
	\maketitle
	\begin{center}
		\includegraphics[scale=0.7]{3.jpg}
	\end{center}

	\newpage
	
	\title{\huge \textit{Representando problemas con listas }}\\
	
	El resumen del día jueves 24 de Enero de 2019 abarca lo siguiente:
	\begin{enumerate}
		\item Conceptos y datos importantes de la clase
			\begin{itemize}
				\item Una forma interesante de darle formato a una lista y que se asemeje a una matriz es usando numpy, y se manda llamar de esta manera: \textit{print (np.matrix(L))}\\
				

				\item Para instalar numpy basta con escribir esto en la terminal: \textit{sudo pip install numpy}\\
				
				\item Cuando estas escribiendo un codigo muy largo en python y no quieres que se vea de esa manera, puedes usar la diagonal invertida y continuar abajo con el código.\\
				
				\item Una variable local es la que existe mientras hago el llamado a la función o mientras se este ejecutando. Y quedaría de la siguiente manera:  \begin{lstlisting}
v0 = 34
g = 9.81
t = 4.3

def posicion (t,v0):
    y = v0*t - 1.0/2*g*t**2
    return(y)

    Donde v0, g y t, son variables locales			
				\end{lstlisting}
				
				
				
			\end{itemize}
		
		
		
		\item Código de python
			\begin{itemize}
				\item \begin{lstlisting}
adn= 'ATGCGACCTAT'
base = 'C'				
def contar_v1(adn,base):
    adn=list(adn)
    i = 0
    for c in adn:
        if c == base:
        i += 1
    return i
				\end{lstlisting} Esta función lo que hace es contar cuántas veces el valor del índice es igual a la base, y se va almacenando en \textit{i} pero previamente convierte la cadena en lista.\\
				
				
				\item \begin{lstlisting}
adn= 'ATGCGACCTAT'
base = 'C'
def contar_v2(adn,base):
    #adn=list(adn)
    i = 0
    for c in adn:
        if c == base:
            i += 1
    return i
				\end{lstlisting} Similar al anterior, pero este no recurre a convertir la cadena en una lista, y esto es porque de igual manera las cadenas usan indices, entonces no fue necesario. Y lo de mas es igual, usa un auxiliar \textit{i} como contador, que vaya almacenando las veces en que la base es igual a algun caracter en la cadena de ADN.\\
				
				\item \begin{lstlisting}
adn= 'ATGCGACCTAT'
base = 'C'
def contar_v3(adn,base):
    #adn=list(adn)
    i = 0
    # for c in adn:
    for j in range(len(adn)):
        if adn[j] == base:
            i += 1
    return i
				\end{lstlisting} De igual manera hace lo mismo que los anteriores, solo que no convierte la cadena en lista y usa la función range, además de la longitud de la cadena, y compara los índices con la base, si es igual entonces suma 1 en i.\\
				
				
				\item \begin{lstlisting}
adn= 'ATGCGACCTAT'
base = 'C'
def contar_v4(adn,base):
   #adn=list(adn)
   i = 0
   j= 0
   # for c in adn:
   #for j in range(len(adn)):
   while j < len(adn):
       if adn[j] == base:
          i += 1
       j +=1
return i
				\end{lstlisting} Una versión más de la función que cuenta cuantas veces aparece una letra en la cadena, a diferencia de los anteriores, usa un \textit{while}, dos variables i y j. La manera en la que funciona es que mientras j sea menor a la longitud de la cadena, el ciclo se va repetir ese total de veces, en el bloque compara si el índice es igual a la base, si es igual almacena +1 en i; la j cuenta cuantas veces se repitio la función, si le ponemos \textit{return i, j} regresa el valor, es decir el total de elementos de la cadena. \\
			
					
				\item \begin{lstlisting}
import numpy as np
print (np.matrix(L)) 
				\end{lstlisting}  Muy útil para darle formato a una lista como si fuera una matriz. y queda de la siguiente manera.\\
				\begin{lstlisting}
       [[ True  True  True  True]
        [False False False  True]
        [ True  True False  True]]
				\end{lstlisting}
				
\newpage
				
				\item \begin{lstlisting}
def resolver(L,e):
   print e
   n = len(L[0])
   m=len(L)
   x = e[0]
   y = e[1]
   #if hay salida:
   if y == n-1 or x == m-1:
       return e[0]+1,e[1]+1 
   else:
       if L[x][y+1] == False:
          e = [x,y+1]
          return resolver(L,e)
  
       elif L[x+1][y] == False:
           e = [x+1,y]
           print e
           return resolver(L,e)
           
       else:
           print "ya no puede avanzar al frente"

type(L)                      
e= [1,0]        
r =resolver(L,e)
				\end{lstlisting} Esta función resuelve el laberinto de 4x3 y que su salida pueda estar abajo o a la derecha. Usa la recursividad para resolver encontrar la salida, partiendo de la coordenada 1,0, que a su vez es una lista e = [1,0]
			
				
				
			\end{itemize}
		
		
		
			
			
		
		
	\end{enumerate}
	
	
	
	
	
	
	
	
	
	
	
\end{document}
